% =============================================================================
% TCC — Algoritmo Genético Híbrido para o CVRPTW em Logística Urbana de Última Milha
% Universidade Paulista (UNIP) — Ciência da Computação — Campus Brasília
%
% Arquivo principal. Compile com:  latexmk main.tex   (ou pela extensão LaTeX Workshop)
% Cada capítulo fica em capitulos/, os elementos pré-textuais em pretextual/.
% =============================================================================

\documentclass[
	12pt,
	openright,
	oneside,
	a4paper,
	english,
	brazil
]{abntex2}

% -----------------------------------------------------------------------------
% Pacotes
% -----------------------------------------------------------------------------
\usepackage[T1]{fontenc}
\usepackage[utf8]{inputenc}
\usepackage{mathptmx}          % Times (mesma fonte do PDF original)
\usepackage{indentfirst}
\usepackage{color}
\usepackage{graphicx}
\usepackage{microtype}
\usepackage{booktabs}          % regras das tabelas
\usepackage{amsmath}
\usepackage[alf,abnt-emphasize=bf]{abntex2cite}  % citações autor-data (ABNT), títulos em negrito

% -----------------------------------------------------------------------------
% Dados do trabalho
% -----------------------------------------------------------------------------
\titulo{Algoritmo Genético Híbrido para o CVRPTW em Logística Urbana de Última Milha}
\autor{Eduardo Gabriel de Sales Rodrigues \\ Esther Chen \\ Rafael Silva dos Santos}
\local{Brasília --- DF}
\data{2026}
\orientador{Prof. Me. Fabricio Rodrigues Freire}
\instituicao{Universidade Paulista --- UNIP}
\tipotrabalho{Trabalho de Conclusão de Curso}
\preambulo{Trabalho de Curso apresentado para obtenção do título de Bacharel em
Ciência da Computação da Universidade Paulista --- UNIP, Campus Brasília.}

% -----------------------------------------------------------------------------
% Metadados do PDF
% -----------------------------------------------------------------------------
\definecolor{blue}{RGB}{41,5,195}
\makeatletter
\hypersetup{
	pdftitle={\@title},
	pdfauthor={Eduardo Gabriel de Sales Rodrigues, Esther Chen, Rafael Silva dos Santos},
	pdfsubject={\imprimirpreambulo},
	pdfcreator={LaTeX with abntex2},
	pdfkeywords={abnt}{latex}{unip}{ciência da computação}{CVRPTW}{algoritmo genético híbrido}{busca local}{heurística gulosa}{otimização combinatória},
	colorlinks=false,
	pdfborder={0 0 0},
	bookmarksdepth=4
}
\makeatother

% Fontes dos títulos: Times em negrito (padrão do abnTeX2 é sem serifa)
\renewcommand{\ABNTEXchapterfont}{\rmfamily\bfseries}

% Recuo das alíneas a), b), c)...
\setlist[alineas,1]{label={\alph*)},topsep=0pt,itemsep=0pt,leftmargin=29pt}

% Espaçamento entre as assinaturas da banca
\setlength{\ABNTEXsignskip}{1.92cm}

% Espaçamento após o indicador de seção e parágrafos
\setlength{\parindent}{1.3cm}
\setlength{\parskip}{0cm}

\makeindex

% =============================================================================
\begin{document}

\selectlanguage{brazil}
\frenchspacing

% ----------------------------------------------------------
% Elementos pré-textuais
% ----------------------------------------------------------
% -----------------------------------------------------------------------------
% Capa (modelo UNIP)
% -----------------------------------------------------------------------------
\begin{capa}
	\centering

	\includegraphics[width=0.5\textwidth]{figuras/estaticas/logo-unip}

	\vspace{0.24cm}

	{\small
	UNIVERSIDADE PAULISTA --- UNIP \\[0.1cm]
	INSTITUTO DE CIÊNCIAS EXATAS \\[0.1cm]
	CURSO DE CIÊNCIA DA COMPUTAÇÃO \par}

	\vspace{1.85cm}

	{\large\bfseries \imprimirtitulo \par}

	\vspace{1.05cm}

	{\small\itshape \imprimirtipotrabalho \par}

	\vspace{2cm}

	{\setlength{\baselineskip}{23pt}\imprimirautor\par}

	\vfill

	{\small\bfseries BRASÍLIA --- DF \par}
	{\small \imprimirdata \par}

	\vspace*{0.75cm}
\end{capa}

\setcounter{page}{2} % a capa conta como página 1
% -----------------------------------------------------------------------------
% Folha de rosto / aprovação (modelo UNIP)
% -----------------------------------------------------------------------------
\begin{folhadeaprovacao}
	\begin{center}
		{\setlength{\baselineskip}{18pt}\imprimirautor\par}

		\vspace*{1cm}

		{\large\bfseries \imprimirtitulo \par}
	\end{center}

	\vspace*{0.1cm}

	% Preâmbulo alinhado à direita, como no modelo da UNIP
	\noindent\hspace*{.466\textwidth}%
	\makebox[0pt][l]{%
		\begin{minipage}[t]{.58\textwidth}
			\small \SingleSpacing
			\imprimirpreambulo

			\vspace{0.4cm}
			Orientador: \imprimirorientador
		\end{minipage}%
	}

	\vspace*{0.55cm}

	\begin{center}
		\small
		Aprovado em: \rule{0.6cm}{0.4pt} de \rule{3.5cm}{0.4pt} de \rule{1.2cm}{0.4pt}

		\vspace*{1.5cm}

		\textbf{BANCA EXAMINADORA}
	\end{center}

	\small
	\vspace*{-1.2cm}
	\assinatura{\imprimirorientador{} --- Orientador \\ Universidade Paulista --- UNIP}
	\assinatura{Prof. Me. \rule{6cm}{0.4pt} \\ Universidade Paulista --- UNIP}
	\assinatura{Prof. Esp. \rule{6cm}{0.4pt} \\ Universidade Paulista --- UNIP}
\end{folhadeaprovacao}

% -----------------------------------------------------------------------------
% Resumo em português
% -----------------------------------------------------------------------------
\begin{resumo}
	Resumo em português
\end{resumo}

% -----------------------------------------------------------------------------
% Abstract (resumo em inglês)
% -----------------------------------------------------------------------------
\begin{resumo}[Abstract]
	\begin{otherlanguage*}{english}
		Resumo em inglês
	\end{otherlanguage*}
\end{resumo}


\pdfbookmark[0]{\contentsname}{toc}
\tableofcontents*
\cleardoublepage

% ----------------------------------------------------------
% Elementos textuais
% ----------------------------------------------------------
\textual

% =============================================================================
% Capítulo 1 — Introdução
% =============================================================================
\chapter{INTRODUÇÃO}
\label{cap:introducao}

Aguarda texto

% =============================================================================
% Capítulo 2 — Fundamentação Teórica
% =============================================================================
\chapter{FUNDAMENTAÇÃO TEÓRICA}
\label{cap:fundamentacao}

Aguarda texto

% =============================================================================
% Capítulo 3 — Metodologia
% Estruturado segundo as etapas de avaliação de desempenho de Jain (1991).
% =============================================================================
\chapter{METODOLOGIA}
\label{cap:metodologia}

% -----------------------------------------------------------------------------
\section{Delineamento da Pesquisa}
\label{sec:delineamento}

Esta pesquisa adota o \textbf{delineamento experimental computacional}, estruturado
segundo o processo sistemático de avaliação de desempenho proposto por
\citeonline{jain1991} em sete etapas sequenciais: definição do sistema e objetivos;
seleção de métricas; identificação de fatores e níveis; escolha da técnica de
avaliação; desenho do experimento; execução e análise dos dados; e apresentação dos
resultados. Cada etapa corresponde a uma seção própria deste capítulo, conferindo
rastreabilidade e reprodutibilidade ao trabalho.

A pesquisa desenvolve-se em duas fases complementares. A \textbf{primeira fase}, de
natureza bibliográfica, compreende o levantamento e a análise das referências que
fundamentam as decisões de modelagem: teoria de otimização combinatória, o Problema de
Roteamento de Veículos com Janelas de Tempo (\textit{CVRPTW}) e o estado da arte em
meta-heurísticas evolutivas aplicadas à logística urbana. A \textbf{segunda fase}, de
natureza experimental, submete os algoritmos propostos a um protocolo de
\textit{benchmarking} sobre instâncias sintéticas calibradas, gerando os dados
quantitativos que sustentam a avaliação da hipótese central.

% -----------------------------------------------------------------------------
\section{Natureza e Classificação da Pesquisa}
\label{sec:natureza}

Quanto à natureza, a pesquisa classifica-se como \textbf{aplicada} \cite{gil2022}: seu
produto é um motor de otimização de rotas concebido para sistemas de logística urbana
com restrições de capacidade e janelas de tempo, com potencial direto de implantação
em operações de entrega de última milha. A contribuição técnico-científica reside na
adaptação formal do \textit{Capacitated Vehicle Routing Problem with Time Windows}
(CVRPTW) ao contexto de entregas urbanas com clusters de demanda, incorporando à
avaliação comparativa duas dimensões operacionais: a taxa de entregas não atendidas e
o custo total de deslocamento da frota.

Quanto aos objetivos, a pesquisa é \textbf{experimental}: define variáveis
independentes controladas pelo experimentador --- os algoritmos de roteamento e as
configurações de parâmetros do Algoritmo Genético Híbrido ---, variáveis dependentes
mensuráveis --- as métricas de eficiência e viabilidade computacional --- e variáveis
controladas fixadas em todas as execuções para garantir a validade interna do
experimento.

Quanto à abordagem, a pesquisa é \textbf{quantitativa}: todos os dados produzidos
pelas execuções são numéricos e submetidos a tratamento estatístico formal, incluindo
estatísticas descritivas, intervalos de confiança, testes de hipótese
não-paramétricos e análise fatorial completa dos parâmetros do Algoritmo Genético,
conforme prescrito por \citeonline{jain1991}.

% -----------------------------------------------------------------------------
\section{Etapa 1 - Definição do Sistema e Objetivos}
\label{sec:etapa1}

O sistema definido consiste em um motor computacional de otimização de rotas para o
Problema de Roteamento de Veículos Capacitados com Janelas de Tempo (CVRPTW),
desenvolvido para apoiar operações de logística urbana de última milha. O sistema
processa instâncias do problema contendo as informações necessárias para o
planejamento das rotas e gera soluções que respeitam as restrições de capacidade dos
veículos e das janelas de tempo dos clientes.

Os objetivos concentram-se na comparação busca verificar se a abordagem híbrida é
capaz de produzir soluções de melhor qualidade para o CVRPTW, reduzindo custos
operacionais e aumentando a eficiência do atendimento sem comprometer o desempenho
computacional exigido para aplicações práticas de logística urbana.

Objetivos da avaliação

A avaliação experimental é conduzida a partir de duas dimensões de desempenho:

Os dois objetivos do experimento são:

\begin{alineas}
	\item \textbf{Objetivo de eficiência}: Avaliar se o Algoritmo Genético Híbrido
	(GA+LS) supera a heurística do Vizinho Mais Próximo, utilizada como baseline, na
	resolução do CVRPTW. A eficiência será medida por meio do custo total de
	deslocamento da frota e da taxa de entregas não atendidas, buscando identificar a
	capacidade do algoritmo híbrido de produzir soluções com menor custo operacional e
	maior nível de atendimento aos clientes.

	\item \textbf{Objetivo de viabilidade computacional}: A viabilidade computacional do
	Algoritmo Genético Híbrido será avaliada a partir do tempo de execução necessário
	para obtenção da solução pelo algoritmo. O tempo de execução será medido em todas
	as instâncias experimentais, considerando múltiplas execuções independentes para
	obtenção da média, variabilidade e pior caso observado. Como critério de
	viabilidade, será adotado um limite máximo de 1500 ms por instância, considerando
	que sistemas de planejamento de rotas urbanas devem fornecer respostas em tempo
	reduzido para permitir a geração ou atualização das rotas sem comprometer o
	processo operacional.
\end{alineas}

% -----------------------------------------------------------------------------
\section{Etapa 2 - Seleção das métricas}
\label{sec:etapa2}

De acordo com \citeonline{jain1991}, a seleção das métricas consiste em definir os
critérios que permitem avaliar o desempenho do sistema frente aos objetivos
estabelecidos. Neste trabalho, essas métricas são utilizadas para comparar o motor de
otimização proposto, baseado no Algoritmo Genético Híbrido (GA+LS), com a heurística
gulosa do Vizinho Mais Próximo, adotada como baseline.

Das instâncias processadas por cada algoritmo --- compostas pelos dados de entrada
(coordenadas, demandas, janelas de tempo e frota), pelo processamento realizado até a
construção da rota e pela saída correspondente ao conjunto de rotas gerado --- são
extraídas quatro métricas: o custo total de deslocamento e a taxa de entregas não
atendidas, obtidos a partir da saída de cada execução, e o tempo de execução e o
desvio das janelas de tempo, associados ao esforço de processamento necessário para
produzir a solução.

Cada uma dessas métricas é detalhada a seguir, contemplando sua definição, o critério
de comparação entre os algoritmos e o procedimento utilizado para sua obtenção.

\begin{alineas}
	\item \textbf{Custo total de deslocamento da frotal}: O custo total de deslocamento
	corresponde à distância acumulada percorrida pela frota na execução das rotas
	geradas, traduzindo diretamente o objetivo de minimização do custo operacional do
	problema de roteamento. O desempenho será comparado considerando que menores
	valores de deslocamento indicam soluções mais eficientes. Para sua obtenção,
	somam-se as distâncias entre os nós visitados em sequência por cada veículo, com
	base nas coordenadas geográficas da rota, até o retorno ao depósito, apurando-se o
	valor total ao final de cada execução.

	\item \textbf{Taxa de entregas não atendidas}: A taxa de entregas não atendidas mede
	a proporção de clientes que não puderam ser incluídos em nenhuma rota válida,
	expondo a capacidade do algoritmo de produzir soluções viáveis sob restrições
	rígidas de capacidade e de janela de tempo. O desempenho será considerado superior
	quando o algoritmo apresentar menores taxas de entregas não atendidas. Sua obtenção
	consiste em contar, ao final de cada execução, os clientes não alocados a nenhuma
	rota viável e dividir esse número pelo total de clientes da instância.

	\item \textbf{Tempo de execução computacional}: O tempo de execução corresponde ao
	intervalo necessário para que o algoritmo produza uma solução completa para uma
	instância, atestando a viabilidade computacional da abordagem para aplicações reais
	de planejamento de rotas urbanas. A avaliação considera múltiplas execuções
	independentes, das quais se extraem o tempo médio, a variabilidade e o pior caso
	observado, sendo o algoritmo viável quando o tempo obtido for inferior ao limite de
	1500 ms por instância. Sua medição é feita cronometrando o intervalo entre o
	recebimento da instância e a produção da solução final, registrado em
	milissegundos a cada execução. Os três valores obtidos são comparados entre a
	heurística do Vizinho Mais Próximo e o Algoritmo Genético Híbrido (GA+LS),
	compondo a base da avaliação experimental apresentada nos capítulos seguintes.

	\item \textbf{Desvio das janelas de tempo}: O desvio das janelas de tempo representa
	o tempo de atraso acumulado nos atendimentos realizados fora do horário limite
	estabelecido, aplicável nos casos em que o modelo permite tolerância de atendimento
	com penalização, em vez de descartar o cliente como não atendido. O desempenho será
	considerado superior quando o algoritmo apresentar menores valores de desvio
	acumulado, indicando maior aderência aos horários programados. Para sua obtenção,
	calcula-se a diferença entre o horário de atendimento efetivamente realizado e o
	horário limite da janela de cada cliente, somando-se os desvios de todos os
	atendimentos com atraso ao final de cada execução.
\end{alineas}

% -----------------------------------------------------------------------------
\section{Etapa 3 - Fatores e níveis}
\label{sec:etapa3}

De acordo com \citeonline{jain1991}, esta etapa consiste em identificar os fatores que
influenciam o desempenho do sistema avaliado e classificá-los em dois grupos: os
parâmetros, fatores cujos níveis são deliberadamente variados ao longo do experimento
com o propósito de observar seu efeito sobre as métricas definidas na Etapa 2, e os
fatores fixos, mantidos constantes em todas as execuções, de modo a evitar que sua
variação introduza ruído estatístico e comprometa a validade interna do experimento.

\begin{alineas}
	\item \textbf{Fator primário: tipo de algoritmo de roteamento}

	O fator central deste experimento é o algoritmo empregado na geração das rotas,
	apresentando dois níveis: a heurística do Vizinho Mais Próximo, adotada como
	baseline, e o Algoritmo Genético Híbrido (GA+LS), proposta principal deste
	trabalho. Este fator sustenta diretamente o objetivo de eficiência definido na
	Etapa 1 (item~\ref{sec:etapa1}), sendo sua variação a que fundamenta a avaliação da
	hipótese central da pesquisa.

	\item \textbf{Configuração do Algoritmo Genético Híbrido}

	Os seguintes parâmetros de configuração do Algoritmo Genético Híbrido são fixados
	em um valor único para todas as execuções deste trabalho, reservando a investigação
	sistemática de seus efeitos sobre o desempenho para trabalhos futuros (item 8):

	\begin{itemize}
		\item Tamanho da população: 100 indivíduos;
		\item Número de gerações: 150;
		\item Taxa de crossover: 0,8;
		\item Taxa de mutação: 0,1;
		\item Peso da penalidade de desvio de janela de tempo na função objetivo: 50;
		\item Uso da busca local (2-opt): ativa.
	\end{itemize}

	Para o fator ``uso da busca local'', contudo, adota-se um nível adicional de
	controle: a desativação completa do componente de busca local, reduzindo o
	algoritmo híbrido a um Algoritmo Genético puro. A comparação entre os dois níveis
	--- busca local ativa (GA+LS) e desativada (GA) --- permite isolar estatisticamente
	a contribuição da hibridização sobre o desempenho, fortalecendo a validade da
	conclusão obtida quanto à eficácia da abordagem híbrida proposta.

	\item \textbf{Fatores relativos à instância do problema}

	Constitui fator adicional a instância do benchmark de Solomon (1987) sobre a qual os
	algoritmos são executados, pertencente ao grupo C1, caracterizado por clientes com
	distribuição espacial agrupada e janelas de tempo amplas. A capacidade do veículo e
	o número de clientes atendidos são determinados pela especificação oficial da
	instância selecionada, não configurando fatores manipulados de forma independente.

	\item \textbf{Fatores fixos}

	Os seguintes fatores são mantidos constantes em todas as execuções do experimento:

	\begin{itemize}
		\item Linguagem de implementação (Python) e biblioteca de otimização evolutiva
		utilizada (DEAP);
		\item Ambiente de hardware sobre o qual as execuções são realizadas;
		\item Número de execuções independentes por configuração testada, fixado em 30,
		conforme necessário à obtenção de média, desvio padrão e aplicação dos testes de
		hipótese não-paramétricos previstos no item~\ref{sec:natureza};
		\item Critério de parada do Algoritmo Genético, definido como número fixo de
		gerações, exceto nas execuções em que este parâmetro constitui o próprio fator
		sob análise.
	\end{itemize}

	A semente de geração de números aleatórios (\textit{seed}) constitui exceção
	deliberada entre os fatores fixos: em vez de mantida constante, é sistematicamente
	variada entre as execuções independentes de uma mesma configuração. Essa variação
	garante que cada repetição represente uma amostra estatisticamente independente do
	comportamento estocástico do algoritmo, condição necessária à validade dos cálculos
	de variabilidade e dos testes de hipótese não-paramétricos previstos na
	classificação da pesquisa.
\end{alineas}

% -----------------------------------------------------------------------------
\section{Etapa 4 - Técnicas de avaliação}
\label{sec:etapa4}

\citeonline{jain1991} descreve três técnicas possíveis para a avaliação de desempenho
de sistemas computacionais: modelagem analítica, simulação e medição direta. Este
trabalho descarta a modelagem analítica como técnica principal, uma vez que o Problema
de Roteamento de Veículos Capacitados com Janelas de Tempo pertence à classe dos
problemas NP-difíceis, para os quais não existe expressão fechada capaz de prever, em
tempo computacionalmente viável, a qualidade das soluções produzidas pelas heurísticas
e meta-heurísticas avaliadas. Ainda assim, a modelagem analítica fundamenta as
equações de custo, as restrições de capacidade e janela de tempo, e a função objetivo
do Algoritmo Genético Híbrido, já apresentadas nas seções anteriores deste capítulo.

Adota-se, portanto, a combinação de simulação e medição direta. A simulação consiste na
execução dos algoritmos sobre instâncias sintéticas e padronizadas do benchmark de
Solomon (1987), que reproduzem de forma controlada padrões realistas de distribuição
espacial de clientes, demandas e janelas de tempo observados em operações de logística
urbana, sem depender da coleta de dados operacionais de uma empresa específica. A
medição direta consiste na obtenção das quatro métricas definidas na Etapa 2 --- custo
total de deslocamento, taxa de entregas não atendidas, tempo de execução computacional
e desvio das janelas de tempo --- diretamente da execução do software implementado, em
vez de estimadas teoricamente.

Complementarmente, os resultados obtidos são confrontados com as melhores soluções
conhecidas publicadas na literatura para a instância C101 do mesmo benchmark,
amplamente adotado em estudos de VRPTW. Essa comparação externa situa o desempenho da
proposta não apenas relativamente ao baseline adotado, mas também frente ao estado da
arte da área de roteamento de veículos.

% -----------------------------------------------------------------------------
\section{Etapa 5 - Instâncias dos testes}
\label{sec:etapa5}

A condução dos experimentos deste trabalho utiliza a instância C101 do benchmark de
Solomon (1987), pertencente ao grupo C1 definido como fator na Etapa 3
(item~\ref{sec:etapa3}), caracterizado por clientes com distribuição espacial agrupada
e janelas de tempo amplas. A \autoref{tab:instancias} resume as características da
instância utilizada.

\begin{table}[htb]
	\centering
	\caption{Instância do benchmark de Solomon (1987) utilizada nos testes}
	\label{tab:instancias}
	\begin{tabular}{ccccl}
		\toprule
		Instância & Clientes & Veículos & Capacidade & Padrão de distribuição \\
		\midrule
		C101      & 100      & 25       & 200        & Agrupado, janelas amplas \\
		\bottomrule
	\end{tabular}
\end{table}

A instância C101 assume, adicionalmente, o papel de referência para a comparação
externa com a literatura descrita na Etapa 4 (item~\ref{sec:etapa4}), por possuir
solução ótima amplamente documentada e citada em estudos de VRPTW.

% -----------------------------------------------------------------------------
\section{Etapa 6 - Desenho do experimento}
\label{sec:etapa6}

% TODO: escrever

% -----------------------------------------------------------------------------
\section{Etapa 7 - Análise dos dados}
\label{sec:etapa7}

% TODO: escrever

% -----------------------------------------------------------------------------
\section{Critérios e Avaliação da Hipótese}
\label{sec:criterios}

% TODO: escrever

% -----------------------------------------------------------------------------
\section{Tecnologias utilizadas}
\label{sec:tecnologias}

% TODO: escrever

% =============================================================================
% Capítulo 4 — Modelagem do Algoritmo Genético Híbrido
% =============================================================================
\chapter{MODELAGEM DO ALGORITMO GENÉTICO HÍBRIDO}
\label{cap:modelagem}

Este capítulo formaliza o algoritmo proposto neste trabalho. Inicialmente, apresenta-se
a formulação matemática do Problema de Roteamento de Veículos Capacitados com Janelas de
Tempo, com a função de custo e as restrições de capacidade e de janela de tempo. Em
seguida, descrevem-se os componentes do Algoritmo Genético Híbrido: a representação dos
indivíduos e sua decodificação em rotas, a função de aptidão empregada durante a busca,
os operadores genéticos, a busca local 2-opt e o procedimento de viabilização que produz
a solução final reportada. Por fim, apresenta-se o fluxo completo do algoritmo e a
diferença entre as configurações GA e GA+LS definidas na Etapa 3
(item~\ref{sec:etapa3}).

% -----------------------------------------------------------------------------
\section{Formulação do problema}
\label{sec:formulacao}

O problema é definido sobre um grafo completo $G = (V, A)$, em que $V = \{0, 1, \dots,
n\}$ é o conjunto de vértices e $A = \{(i, j) : i, j \in V, i \neq j\}$ é o conjunto de
arcos. O vértice $0$ representa o depósito, e o conjunto $C = V \setminus \{0\}$
representa os $n$ clientes. A cada arco $(i, j)$ associam-se a distância $d_{ij}$,
calculada pela distância euclidiana entre as coordenadas dos vértices, e o tempo de
viagem $t_{ij}$, que, conforme a convenção do benchmark de \citeonline{solomon1987}, é
numericamente igual à distância. A cada cliente $i \in C$ associam-se a demanda $q_i$, o
tempo de serviço $s_i$ e a janela de tempo $[e_i, l_i]$, em que $e_i$ é o instante mais
cedo e $l_i$ o instante mais tarde para o início do atendimento. A janela do depósito,
$[e_0, l_0]$, corresponde ao horizonte de planejamento. A frota é composta por um
conjunto $K$ de $m$ veículos idênticos, cada um com capacidade $Q$.

Adota-se a formulação clássica de três índices descrita por \citeonline{cordeau2002}. A
variável binária $x_{ijk}$ assume valor 1 quando o veículo $k$ percorre o arco $(i, j)$
e 0 caso contrário, e a variável contínua $w_{ik}$ indica o instante de início do
atendimento do vértice $i$ pelo veículo $k$. O problema consiste em:

\begin{align}
	\min \quad & \sum_{k \in K} \sum_{(i,j) \in A} d_{ij} \, x_{ijk}
	\label{eq:objetivo} \\
	\text{sujeito a} \quad
	& \sum_{k \in K} \sum_{j \in V} x_{ijk} = 1,
	&& \forall i \in C \label{eq:visita} \\
	& \sum_{j \in C} x_{0jk} \leq 1,
	&& \forall k \in K \label{eq:saida} \\
	& \sum_{i \in V} x_{ihk} - \sum_{j \in V} x_{hjk} = 0,
	&& \forall h \in V, \forall k \in K \label{eq:fluxo} \\
	& \sum_{i \in C} q_i \sum_{j \in V} x_{ijk} \leq Q,
	&& \forall k \in K \label{eq:capacidade} \\
	& w_{ik} + s_i + t_{ij} - w_{jk} \leq M (1 - x_{ijk}),
	&& \forall (i, j) \in A, j \neq 0, \forall k \in K \label{eq:tempo} \\
	& e_i \leq w_{ik} \leq l_i,
	&& \forall i \in V, \forall k \in K \label{eq:janela} \\
	& x_{ijk} \in \{0, 1\},
	&& \forall (i, j) \in A, \forall k \in K \label{eq:binaria}
\end{align}

A função objetivo \eqref{eq:objetivo} minimiza a distância total percorrida pela frota,
correspondente à métrica de custo total de deslocamento definida na Etapa 2
(item~\ref{sec:etapa2}). A restrição \eqref{eq:visita} garante que cada cliente seja
visitado exatamente uma vez; a \eqref{eq:saida} limita a uma única rota por veículo,
partindo do depósito; e a \eqref{eq:fluxo} assegura a conservação de fluxo, de modo que
todo veículo que chega a um vértice também o deixa. A restrição \eqref{eq:capacidade}
impõe que a demanda total atendida por um veículo não exceda sua capacidade. A
\eqref{eq:tempo} encadeia os instantes de atendimento ao longo da rota, sendo $M$ uma
constante suficientemente grande que desativa a restrição quando o arco não é
percorrido, e a \eqref{eq:janela} exige que todo atendimento se inicie dentro da janela
do vértice. Por fim, a \eqref{eq:binaria} define o domínio das variáveis de decisão.

Dada uma rota $r = (0, v_1, v_2, \dots, v_p, 0)$, os instantes de atendimento são
obtidos por simulação ao longo da sequência. O instante de chegada ao cliente $v_j$ é
dado por
\begin{equation}
	a_{v_j} = b_{v_{j-1}} + s_{v_{j-1}} + t_{v_{j-1} v_j},
	\label{eq:chegada}
\end{equation}
com $b_{v_0} = e_0$ e $s_{v_0} = 0$ para o depósito, e o início do atendimento é
\begin{equation}
	b_{v_j} = \max\left(a_{v_j}, e_{v_j}\right),
	\label{eq:inicio}
\end{equation}
isto é, um veículo que chega antes da abertura da janela aguarda até o instante
$e_{v_j}$. Quando $b_{v_j} > l_{v_j}$, o atendimento ocorre com atraso, o que viola a
restrição \eqref{eq:janela}. O atraso do cliente e o atraso acumulado da rota são
definidos, respectivamente, por
\begin{equation}
	\lambda_{v_j} = \max\left(0, \, b_{v_j} - l_{v_j}\right)
	\qquad \text{e} \qquad
	\Lambda(r) = \sum_{j=1}^{p} \lambda_{v_j},
	\label{eq:atraso}
\end{equation}
e a distância da rota, por
\begin{equation}
	D(r) = d_{0 v_1} + \sum_{j=1}^{p-1} d_{v_j v_{j+1}} + d_{v_p 0}.
	\label{eq:distancia}
\end{equation}

Nas instâncias de \citeonline{solomon1987}, as janelas dos clientes são definidas de modo
que um veículo que inicia o atendimento até $l_i$ ainda consegue retornar ao depósito
dentro do horizonte $l_0$. Assim, o respeito às janelas dos clientes garante também o
respeito à janela do depósito, e o atraso $\Lambda(r)$ considera apenas os clientes.
% TODO(M1): conferir no parser que l_i + s_i + t_{i0} <= l_0 vale para todo cliente da C101.

% -----------------------------------------------------------------------------
\section{Representação e decodificação}
\label{sec:representacao}

Cada indivíduo do Algoritmo Genético é representado por um cromossomo de permutação
$\pi = (\pi_1, \pi_2, \dots, \pi_n)$, que contém cada cliente de $C$ exatamente uma vez
e não inclui o depósito nem delimitadores de rota. Essa representação, denominada
\textit{giant tour}, foi proposta por \citeonline{prins2004} para o problema de
roteamento de veículos e corresponde a uma única sequência de visitas, como se um só
veículo de capacidade ilimitada atendesse todos os clientes. Sua principal vantagem é
permitir o uso direto de operadores genéticos de permutação, que sempre produzem
cromossomos válidos, sem necessidade de reparo.

A conversão do cromossomo em um conjunto de rotas é feita pelo procedimento
\textit{Split} \cite{prins2004}, que particiona a sequência $\pi$ em subsequências
consecutivas, cada uma atendida por um veículo, preservando a ordem do cromossomo. O
procedimento constrói um grafo auxiliar acíclico $H$ com os vértices $0, 1, \dots, n$,
no qual cada arco $(i, j)$, com $i < j$, representa a rota $r_{ij} = (0, \pi_{i+1},
\dots, \pi_j, 0)$. O arco existe somente se a demanda dessa rota respeita a capacidade
do veículo, isto é, se
\begin{equation}
	\sum_{h = i+1}^{j} q_{\pi_h} \leq Q,
	\label{eq:split-capacidade}
\end{equation}
e seu peso é a contribuição da rota para a função de aptidão, definida na
\autoref{sec:aptidao}. O caminho mínimo do vértice $0$ ao vértice $n$ em $H$,
calculado por programação dinâmica em ordem topológica, corresponde à melhor partição
possível da sequência $\pi$ em rotas, e cada arco desse caminho define uma rota da
solução decodificada.

Como os arcos que excedem a capacidade não existem em $H$, nenhuma solução decodificada
contém rota com demanda superior a $Q$. A restrição de capacidade \eqref{eq:capacidade}
é, portanto, satisfeita por construção, conforme estabelecido na Etapa 2
(item~\ref{sec:etapa2}), e não precisa ser penalizada na função de aptidão. Caso o
caminho mínimo resulte em mais de $m$ rotas, emprega-se a variante do procedimento com
frota limitada, também descrita por \citeonline{prins2004}, que restringe o número de
arcos do caminho a $m$.

% -----------------------------------------------------------------------------
\section{Função de aptidão}
\label{sec:aptidao}

Conforme o tratamento das janelas de tempo definido na Etapa 2 (item~\ref{sec:etapa2}),
o Algoritmo Genético admite, durante a busca, soluções com atendimentos após o limite da
janela, penalizando o atraso na função de aptidão. Para um indivíduo $\pi$ decodificado
no conjunto de rotas $R(\pi)$, a aptidão é dada por
\begin{equation}
	f(\pi) = \sum_{r \in R(\pi)} \left[ D(r) + \beta \, \Lambda(r) \right],
	\label{eq:aptidao}
\end{equation}
em que $D(r)$ é a distância da rota \eqref{eq:distancia}, $\Lambda(r)$ é seu atraso
acumulado \eqref{eq:atraso} e $\beta$ é o peso da penalidade de desvio de janela de
tempo, fixado em 50 na Etapa 3 (item~\ref{sec:etapa3}). O problema é de minimização:
quanto menor o valor de $f$, mais apto é o indivíduo.

Como a aptidão é a soma das contribuições de cada rota, o peso de cada arco $(i, j)$ do
grafo auxiliar do \textit{Split} é $D(r_{ij}) + \beta \, \Lambda(r_{ij})$. Desse modo, a
decodificação escolhe, para a ordem de visitas fixada pelo cromossomo, a partição em
rotas de menor aptidão, e não apenas a de menor distância. Uma solução sem atraso tem
aptidão igual à sua distância total, que coincide com a função objetivo
\eqref{eq:objetivo}.

% -----------------------------------------------------------------------------
\section{Operadores genéticos}
\label{sec:operadores}

O ciclo evolutivo segue a estrutura geral dos Algoritmos Genéticos descrita por
\citeonline{holland1975} e \citeonline{goldberg1989}: uma população de indivíduos é
avaliada, selecionada, recombinada e mutada ao longo de um número fixo de gerações. Os
parâmetros de população, número de gerações e taxas de cruzamento e mutação são os
fixados na Etapa 3 (item~\ref{sec:etapa3}). Os operadores empregados são descritos a
seguir.

\begin{alineas}
	\item \textbf{População inicial}: A população inicial é composta por 100
	permutações aleatórias dos clientes, geradas a partir da semente da execução. Não se
	inclui na população inicial a solução da heurística do Vizinho Mais Próximo, de modo
	que o desempenho do Algoritmo Genético não dependa do baseline com o qual é
	comparado.

	\item \textbf{Seleção}: Os pais são escolhidos por seleção por torneio
	\cite{goldberg1989}, na qual $k = 3$ indivíduos são sorteados da população e o de
	menor aptidão é selecionado, repetindo-se o processo até completar o número de pais
	necessário. A seleção por torneio depende apenas da ordem relativa das aptidões, e
	não de sua escala, o que a torna adequada a uma função de aptidão cujos valores
	variam fortemente com a penalidade de atraso.

	\item \textbf{Cruzamento}: Os pais são recombinados, com probabilidade 0,8 por par,
	pelo cruzamento de ordem (\textit{Order Crossover}, OX) \cite{davis1985}. O operador
	sorteia dois pontos de corte, copia para o filho o segmento do primeiro pai entre
	esses pontos e preenche as posições restantes com os clientes ausentes, na ordem em
	que aparecem no segundo pai. O resultado é sempre uma permutação válida, que
	preserva a ordem relativa das visitas de ambos os pais.

	\item \textbf{Mutação}: Cada filho sofre mutação com probabilidade 0,1, por
	inversão: sorteiam-se duas posições do cromossomo e inverte-se a ordem dos clientes
	entre elas. A inversão altera a sequência de visitas preservando a maior parte das
	adjacências do cromossomo, o que favorece mudanças graduais na solução.

	\item \textbf{Elitismo}: O melhor indivíduo de cada geração é copiado sem alterações
	para a geração seguinte, garantindo que a melhor solução encontrada não seja perdida
	pela ação dos operadores.

	\item \textbf{Critério de parada}: A evolução termina ao completar 150 gerações,
	conforme o fator fixo definido na Etapa 3 (item~\ref{sec:etapa3}).
\end{alineas}

% -----------------------------------------------------------------------------
\section{Busca local 2-opt}
\label{sec:buscalocal}

A hibridização do Algoritmo Genético consiste em aplicar uma busca local aos indivíduos
gerados, abordagem conhecida como algoritmo memético \cite{moscato1989}. Enquanto os
operadores genéticos exploram regiões distintas do espaço de soluções, a busca local
refina cada solução até um ótimo local, compensando a baixa capacidade de ajuste fino
dos operadores de recombinação.

Emprega-se a vizinhança 2-opt \cite{croes1958}, aplicada separadamente a cada rota da
solução decodificada. Um movimento 2-opt remove dois arcos não adjacentes da rota,
$(v_i, v_{i+1})$ e $(v_j, v_{j+1})$, com $i < j$, e os substitui pelos arcos $(v_i,
v_j)$ e $(v_{i+1}, v_{j+1})$, o que equivale a inverter a ordem de visita dos clientes
entre as posições $i+1$ e $j$. Como a inversão altera os instantes de chegada de todos
os clientes posteriores a $v_i$, o efeito do movimento sobre as janelas de tempo não se
reduz à variação de distância. Por isso, cada movimento é avaliado pela contribuição da
rota para a aptidão, $D(r) + \beta \, \Lambda(r)$, e aceito somente quando a reduz.

A busca adota a estratégia de primeira melhoria: o primeiro movimento que reduz a
contribuição da rota é aplicado, e a varredura recomeça a partir da rota modificada,
até que nenhum movimento a melhore. Ao final, as rotas melhoradas são concatenadas, na
ordem em que aparecem na solução, para formar o novo cromossomo do indivíduo, que é
reavaliado pelo \textit{Split}. Como a partição obtida pela busca local é um dos
caminhos possíveis no grafo auxiliar, e o \textit{Split} retorna o de menor peso, a
aptidão reavaliada nunca é pior que a obtida pela busca local.

A busca local é aplicada a todo indivíduo novo: a cada indivíduo da população inicial e
a cada filho modificado por cruzamento ou mutação. Indivíduos copiados sem alteração,
como o de elite, já são ótimos locais e não são submetidos novamente ao 2-opt.

% -----------------------------------------------------------------------------
\section{Viabilização da solução final}
\label{sec:viabilizacao}

Ao término das gerações, o melhor indivíduo encontrado ao longo de toda a execução é
decodificado, e o atraso acumulado de suas rotas, $\sum_{r} \Lambda(r)$, é registrado
como a métrica de desvio das janelas de tempo definida na Etapa 2
(item~\ref{sec:etapa2}). Como, na solução final reportada, as janelas de tempo são
tratadas como restrições rígidas, essa solução passa por um procedimento de
viabilização antes do cálculo das demais métricas.

O procedimento percorre cada rota na ordem de visita, recalculando os instantes de
chegada pelas equações \eqref{eq:chegada} e \eqref{eq:inicio}. Sempre que o instante de
chegada a um cliente excede o limite de sua janela, o cliente é retirado da rota e
contabilizado como não atendido, e o percurso prossegue do último cliente atendido
diretamente para o seguinte. Como as distâncias euclidianas satisfazem a desigualdade
triangular, a retirada de um cliente nunca atrasa a chegada aos clientes posteriores,
de modo que uma única passagem pela rota basta para eliminar todos os atrasos. A
retirada também reduz a carga do veículo, preservando a restrição de capacidade. Rotas
que ficam vazias são descartadas.

A solução resultante não contém atrasos e, sobre ela, calculam-se o custo total de
deslocamento, pela soma das distâncias \eqref{eq:distancia} das rotas restantes, e a
taxa de entregas não atendidas, pela razão entre o número de clientes retirados e o
total de clientes $n$. Os clientes retirados não são reinseridos em outras rotas, o que
torna a taxa de não atendidos uma medida direta da capacidade da busca de produzir
soluções que respeitam as janelas de tempo.

% -----------------------------------------------------------------------------
\section{Fluxo do algoritmo}
\label{sec:fluxo}

O \autoref{quadro:gals} resume o fluxo completo do Algoritmo Genético Híbrido. A
configuração GA, adotada como controle na Etapa 3 (item~\ref{sec:etapa3}), executa o
mesmo fluxo com os passos 3 e 9 desativados. A representação, a decodificação, a
função de aptidão, os operadores genéticos e a viabilização são idênticos nas duas
configurações, de modo que a diferença de desempenho entre elas pode ser atribuída
exclusivamente à busca local.

\begin{quadro}[htb]
	\centering
	\caption{Fluxo do Algoritmo Genético Híbrido (GA+LS)}
	\label{quadro:gals}
	\begin{tabular}{p{0.9\textwidth}}
		\toprule
		\textbf{Entrada:} instância, semente da execução \\
		\textbf{Saída:} rotas viáveis, clientes não atendidos, desvio das janelas \\
		\midrule
		1. Calcular a matriz de distâncias $d_{ij}$ e de tempos $t_{ij}$. \\
		2. Gerar 100 permutações aleatórias dos clientes. \\
		3. Aplicar o 2-opt a cada indivíduo da população inicial. \\
		4. Avaliar cada indivíduo pelo \textit{Split} e pela aptidão $f$. \\
		5. \textbf{Para} cada uma das 150 gerações: \\
		\quad 6. Copiar o melhor indivíduo para a nova população (elitismo). \\
		\quad 7. Selecionar os pais por torneio ($k = 3$). \\
		\quad 8. Aplicar o cruzamento OX (0,8) e a mutação por inversão (0,1). \\
		\quad 9. Aplicar o 2-opt a cada filho modificado. \\
		\quad 10. Avaliar os filhos e formar a nova população. \\
		11. Decodificar o melhor indivíduo encontrado e registrar o desvio. \\
		12. Viabilizar a solução e calcular custo e taxa de não atendidos. \\
		\bottomrule
	\end{tabular}
\end{quadro}

% =============================================================================
% Capítulo 5 — Implementação e Experimentos
% =============================================================================
\chapter{IMPLEMENTAÇÃO E EXPERIMENTOS}
\label{cap:implementacao}

Aguarda texto

% =============================================================================
% Capítulo 6 — Resultados e Discussão
% =============================================================================
\chapter{RESULTADOS E DISCUSSÃO}
\label{cap:resultados}

Este capítulo apresenta e analisa os dados produzidos pelo experimento descrito no
\autoref{cap:implementacao}.

% =============================================================================
% Capítulo 7 — Conclusão
% =============================================================================
\chapter{CONCLUSÃO}
\label{cap:conclusao}

Este trabalho investigou se um Algoritmo Genético Híbrido é capaz de superar a
heurística do vizinho mais próximo na construção de rotas de entrega urbana com
restrições de capacidade e janelas de tempo rígidas, sem comprometer a viabilidade
computacional do sistema.


% ----------------------------------------------------------
% Elementos pós-textuais
% ----------------------------------------------------------
\postextual

\bibliography{referencias}

\begin{apendicesenv}
% =============================================================================
% Apêndice A — Algoritmos
% =============================================================================
\chapter{ALGORITMOS}
\label{apendice:algoritmos}

\end{apendicesenv}

\end{document}
